\scanpage{8}
Posons
\[
\left\{\begin{aligned}
s_0&=f_0\\
s_1&=\sigma_0.f_0\\
s_2&=\sigma_1s_1\\
&\cdots\\
s_i&=\sigma_{i-1}s_{i-1}\\
&\cdots\\
s_{\deleted{\unsure{n}}n+1}&=\sigma_ns_n
\end{aligned}\right.
\]
Je dis que les $(s_i)_{0\leq i\leq n+1}$ \unsure{forment|forme} une \uline{base Affine} de $E$ fibre par fibre, \inserted{et que $f_i=\text{env. aff. de }(s_0\text{--}s_i)$, $0\leq i\leq n$} \deleted{$\unsure{f_i=s_i}$}, et que \inserted{$\sigma_i s_j=s_j$ si $i>j$;} Supposons prouvé, \inserted{pour $i$ donné $0\leq i<n$, que $\{s_0,\ldots,s_j\}$ \unsure{affinemt libre}, $f_j=\text{enveloppe affine de }\{s_0\text{--}s_j\}$, $0\leq j\leq i$,} \deleted{ceci étant trivial pr $i=0$). (Il faut \unsure{prouver}} enfin $\sigma_k s_j=s_j$ si \deleted{$\unsure{j<i}$}\inserted{$k>j$}, \deleted{\unsure{toujours}} pour $j\leq i$. Prouvons que \inserted{fibre par fibre} $s_{i+1}=\sigma_i s_i$ n'appartient pas à \inserted{l'enveloppe} affine $f_i$ de $\{s_0,\ldots,s_i\}$, et \inserted{plus précisément} que $f_{i+1}$ est l'enveloppe affine de $\{s_0,\ldots,s_i,\inserted{s_{i+1}=}\sigma_i s_i\}$, \deleted{\unsure{Et}} puis \inserted{que $\sigma_k s_{i+1}=s_{i+1}$ si $k>i+1$.} \unsure{D'abord|Évident}, comme $s_i\in f_i\subset f_{i+1}$ et $\sigma_i$ fixe $f_{i+1}$, on a $\inserted{s_{i+1}=}\sigma_i(s_i)\in f_{i+1}$, donc il faut seulement prouver que $\sigma_i(s_i)\notin f_i$. \deleted{Mais si}\inserted{\unsure{plus encore que} sur les fibres, $f_{i+1}=\operatorname{Envaff}(s_0,\ldots,s_{i+1})$. Mais si} on avait $\sigma_i(s_i)\in f_i$, \deleted{comme $\unsure{\sigma_i f_{i-1}=f_{i-1}}$, donc $\sigma_i$ On aurait $\sigma_k f_i=f_i$ pr $k\neq i$, on aurait $\unsure{\sigma_k f}$} \inserted{On aurait}
\[
\begin{aligned}
\sigma_i(f_i)&=\sigma_i\unsure{\operatorname{Env}|\operatorname{Ens}}(f_{\unsure{i-1|i+1}},s_i)\deleted{\unsure{=\operatorname{Envaff}_i(}}\\
&=\operatorname{Env}(\underset{f_{i-1}}{\uline{\sigma_i(f_{i-1})}},\underset{\substack{\uparrow\\f_i}}{\sigma_i s_i})\subset f_i
\end{aligned}
\]
donc $\sigma_i(f_i)=f_i$, donc comme $\sigma_j f_i=f_i$ pour $j\neq i$, on aurait $\sigma_k f_i=f_i$ pour $0\leq k\leq n$, donc $f_i$ fixé par $\Gamma=G$ donc
% archive sheet number: 7
